\chapter{The Memory and Knowledge Layer}

The memory and knowledge layer determines what information is available to the model at
the moment of decision. Its function is to construct the context window with the subset
of available information that is most relevant and most likely to support a correct and
grounded response. Context engineering is the discipline of designing, assembling, and
maintaining the information environment available to an agent at inference time, curating
the optimal set of tokens during inference, including information beyond the prompt
itself.

\section{Why Memory and Knowledge Form a Single Layer}

Memory and knowledge are distinct in origin but unified in function. Memory accumulates
through interaction: it records user preferences, past decisions, prior conversation
turns, and outcomes. Knowledge is sourced externally: documents, policies, codebases,
manuals and databases. Both must pass through the same context engineering
process before reaching the model's context window. Context engineering governs this
selection, deciding which pieces of memory and which fragments of knowledge to include
for a given task, in what form, and at what level of compression.

The two are treated as a single layer because failures in either produce the same
downstream effect on grounding: the model reasons from an incomplete or misleading
context window. An agent that retrieves irrelevant document chunks, omits a critical
policy, or fails to surface a previously stated user preference will produce ungrounded
outputs regardless of the model's underlying capability.

\section{Types of Memory}

Memory in agentic systems is usefully classified by its temporal scope and function.
Working memory is the immediate task context: the current messages, active plan, tool
outputs, and unresolved constraints that the agent is reasoning over in the present turn.
Short-term memory spans a session: it preserves conversational history across turns and
enables the agent to maintain coherence without re-establishing context on every exchange.
Episodic memory records past events across sessions: what the user requested, what the
agent did, and what outcome resulted. Semantic memory holds stable extracted facts: user
role, project names, known issues, and domain definitions that remain valid across
interactions. Preference memory captures standing instructions about tone, format, default
tools, and preferred workflows. Procedural memory contains reusable packages of know-how:
skills, checklists, playbooks, and workflow templates that the agent can retrieve and
apply when a task matches their scope.

\section{Skills as Procedural Knowledge}

A skill occupies a specific position in this taxonomy. It is neither a model weight nor
a tool: it is a retrievable unit of procedural knowledge that the system injects into the
context window when the current task requires it. A skill may take the form of a
domain-specific instruction fragment, a step-by-step checklist for a recurring workflow,
a repository development guide, a troubleshooting playbook, or a bundle of tool-use rules
paired with expected output formats.

The key property of a skill is that it is selected, versioned, and evaluated independently
of the model. A skill whose content is outdated introduces incorrect procedural guidance,
and one that is irrelevant to the current task adds noise to the context window and
degrades grounding. Skills therefore belong in the knowledge layer and are subject to the
same retrieval, quality, and freshness requirements as any other document.

A common representation is a structured Markdown document that pairs machine-readable
frontmatter with human-readable procedural content. The frontmatter declares the skill's
name, its intended scope, and the tools it is permitted to invoke; the body encodes the
step-by-step procedure the agent should follow when the skill is active:

\begin{lstlisting}[language={}]
---
name: data-pipeline-diagnosis
description: >
  Procedure for diagnosing failures in a data processing pipeline.
  Use when the user reports missing, delayed, or malformed output.
allowed-tools:
  - fetch_job_logs
  - get_job_status
  - inspect_schema
---

## Diagnosis Flow

1. Retrieve the job identifier from the user or from recent job history.
2. Call get_job_status to determine whether the job completed,
   failed, or is still running.
3. If the job failed, call fetch_job_logs and identify the first
   error entry.
4. Call inspect_schema to verify that the input data conforms to
   the expected format.
5. Summarise the root cause and propose a remediation step.
\end{lstlisting}

This format makes the skill inspectable, diffable, and testable. A new version of the
skill can be evaluated against a dataset of past failures before it replaces the previous
version in the retrieval index. The agent framework retrieves the skill document when the
current task matches its declared scope and injects its content into the context window
alongside the user's request.

\section{Context Compression}

Context windows are finite and carry a cost proportional to their length. Compression is
the process of reducing the volume of context while preserving the information most
relevant to the task at hand. Common techniques include summarising prior conversation
turns through memory consolidation, retaining only the highest-ranked retrieval chunks,
and converting verbose tool outputs into compact structured fields containing only the
decision-relevant values.

Compression introduces its own failure modes. A consolidation step may omit a critical
qualification present in the original. A memory extraction step may store a temporary
instruction as a permanent preference. A retrieval filter may exclude the document that
contains the correct grounding evidence. For this reason, memory and context
transformations require the same evaluation discipline as model outputs: they are active
interventions in the context engineering pipeline that can introduce errors at any step.

\section{Memory Write Policy}

Selective retention is a design requirement. An agent that stores every piece of
information it encounters produces a memory that degrades in utility as it accumulates
noise, outdated facts, and low-confidence inferences. A principled write policy stores
information when the user has explicitly requested retention, when the fact is stable and
reusable across future sessions, when storage improves personalisation or session
continuity in a way the user has consented to, and when the information originates from
a trusted source. Sensitive information that serves no persistent purpose, temporary
context that applies only to the current task, and content inferred with low confidence
warrant exclusion from long-term memory.

The system should expose memory to inspection and correction. A user who can query what
the agent has retained and modify or delete individual entries has meaningful control over
the grounding context the agent will assemble in future sessions.

\section{Retrieval-Augmented Generation}

Retrieval-Augmented Generation (RAG) is the standard pattern for grounding model
responses in external knowledge. It proceeds through four operations. Chunking divides
source documents into retrieval units of appropriate granularity, each carrying metadata
that identifies its origin, recency, and scope. Embedding maps each chunk into a vector
representation that supports semantic similarity search. Retrieval selects the candidate
chunks most relevant to the current query, typically by combining vector similarity with
keyword or metadata filters. Context assembly arranges the selected chunks into the
context window in a form the model can reason from, managing ordering, deduplication, and
length constraints.

Each of these operations is a potential source of grounding failure and should be tested
independently. A chunking strategy that splits sentences across unit boundaries degrades
retrieval precision. An embedding model that has drifted from the domain of the documents
produces unreliable similarity scores. A retrieval step that returns the highest-scoring
chunks without diversity may miss the document that contains the correct answer. A context
assembly step that places the most relevant content at the end of a long context window
may cause the model to underweight it due to the well-documented lost-in-the-middle
effect. Treating RAG as a pipeline of testable components rather than a monolithic
retrieval call is what makes its failure modes diagnosable.

\section{Summary}

The memory and knowledge layer is where skills, context engineering, long-term memory,
and retrieval infrastructure belong. Its function is to construct the context window with
the right information at the right time. Production agent design is therefore as much a
context engineering discipline as a model engineering one: the quality of the context
assembled for each task is as consequential as the capability of the model that receives
it.